\subsection{Trees and Ensemble Averaging}

Tree-based methods partition the input space into simple regions and fit a local rule in each region.  Ensemble averaging methods keep this local modelling idea but stabilise it by resampling, voting, and decorrelating the individual trees.

\begin{definition} \label{def:sl-decision-tree} \hlt{Decision Tree}\\
Let $\mathcal{X}\subseteq\R^p$ and let $\mathcal{A}$ be the prediction space.\\
A decision tree $T$ is a finite binary recursive partition of $\mathcal{X}$ with a leaf value in $\mathcal{A}$ on each terminal region.
\begin{enumerate}[label=(\roman*)]
\setlength{\itemsep}{0pt}
\item \hlt{Split}: a split at node $R\subseteq\mathcal{X}$ is a feature-threshold pair $(j,t)$ with $j\in\{1,\ldots,p\}$ and $t\in\R$.\\
It produces children $R_L\defeq R\cap\{x:x_j\leq t\}$ and $R_R\defeq R\cap\{x:x_j>t\}$.
\item \hlt{Terminal partition}: after splitting stops, the leaves of $T$ are nonempty regions $R_1(T),\ldots,R_M(T)$\\
with $\mathcal{X}=\bigcup\limits_{m=1}^M R_m(T)$ and $R_m(T)\cap R_{m'}(T)=\varnothing$ for $m\neq m'$.
\item \hlt{Leaf rule}: a leaf assignment $a_m\in\mathcal{A}$ for $m=1,\ldots,M$ defines $\widehat{f}_T:\mathcal{X}\to\mathcal{A}$ by $\widehat{f}_T(x)\defeq a_m$\\
for the unique $m$ such that $x\in R_m(T)$.
\end{enumerate}
\end{definition}

\begin{definition} \label{def:sl-bagging-random-forest} \hlt{Bagging and Random Forests}\\
Let $D_n=\{(x_i,y_i)\}_{i=1}^n$, $B\geq1$, and $m\leq p$.
\begin{enumerate}[label=(\roman*)]
\setlength{\itemsep}{0pt}
\item \hlt{Bagging}: draw $D_1^*,\ldots,D_B^*$ from $D_n$, fit $T_1,\ldots,T_B$, and aggregate $\widehat{f}_{\mathrm{bag}}(x)\defeq B^{-1}\sum\limits_{b=1}^B T_b(x)$ for regression, with voting or averaged probabilities for classification.
\item \hlt{Random forest}: in addition to bagging, at each node $R$ of $T_b$ choose $M_{b,R}\subseteq\{1,\ldots,p\}$ with $|M_{b,R}|=m$ and optimise the split over features in $M_{b,R}$ only.
\end{enumerate}
Bagging averages bootstrap perturbations; random split-feature subsets reduce $\operatorname{corr}(T_b,T_{b'})$ before averaging.
\end{definition}

\begin{breakablealgorithm}
\caption{Random Forest}
\label{alg:sl-random-forest}
\begin{algorithmic}
\Require Training data $(x_i,y_i)_{i=1}^n$, number of trees $B$, candidate split count $m\leq p$, tree stopping rule
\Ensure Trees $T_1,\ldots,T_B$ and an aggregated predictor
\For{$b=1,\ldots,B$}
\State Draw a bootstrap resample $S^{\ast b}$ of size $n$
\State Grow a tree $T_b$ on $S^{\ast b}$
\For{each nonterminal node while growing $T_b$}
\State Select $m$ of the $p$ predictors at random
\State Choose the best split using only those $m$ predictors
\State Split the node and continue until the stopping rule is met
\EndFor
\EndFor
\State For regression, return $\widehat{f}_{\mathrm{rf}}(x)=B^{-1}\sum_{b=1}^B T_b(x)$
\State For classification, return a majority vote or averaged class-probability rule from $T_1(x),\ldots,T_B(x)$
\end{algorithmic}
\end{breakablealgorithm}

\begin{figure}[H]
\centering
\begin{tikzpicture}[
    >=Stealth,
    font=\scriptsize,
    line cap=round,
    line join=round,
    box/.style={draw=gray!65, rounded corners=2pt, align=center, minimum height=0.5cm, inner xsep=6pt, inner ysep=2pt},
    data/.style={box, fill=gray!8, minimum width=1.85cm},
    sample/.style={box, draw=cyan!65!black, fill=cyan!8, minimum width=2.05cm},
    tree/.style={box, draw=green!45!black, fill=green!8, minimum width=1.45cm},
    agg/.style={box, draw=orange!70!black, fill=orange!10, minimum width=1.65cm},
    oobbox/.style={box, draw=red!60!black, fill=red!6, minimum width=1.75cm},
    flow/.style={->, draw=black!75, line width=0.6pt},
    oob/.style={->, dashed, draw=red!70!black, line width=0.6pt}
]
\node[data] (data) at (0,0) {training set};
\node[sample] (s1) at (2.55,0.85) {bootstrap $S^{\ast 1}$};
\node[sample] (s2) at (2.55,0) {bootstrap $S^{\ast 2}$};
\node[sample] (s3) at (2.55,-0.85) {bootstrap $S^{\ast 3}$};

\node[tree] (t1) at (5.05,0.85) {tree $T_1$};
\node[tree] (t2) at (5.05,0) {tree $T_2$};
\node[tree] (t3) at (5.05,-0.85) {tree $T_3$};
\node[agg] (agg) at (7.35,0) {aggregate};

\node[oobbox] (held) at (3.15,-1.6) {OOB obs. $i$};
\node[oobbox] (oobavg) at (7.35,-1.6) {OOB prediction};

\draw[flow] (data.north east) -- (s1.west);
\draw[flow] (data.east) -- (s2.west);
\draw[flow] (data.south east) -- (s3.west);
\draw[flow] (s1.east) -- (t1.west);
\draw[flow] (s2.east) -- (t2.west);
\draw[flow] (s3.east) -- (t3.west);
\draw[flow] (t1.east) -- (agg.north west);
\draw[flow] (t2.east) -- (agg.west);
\draw[flow] (t3.east) -- (agg.south west);

\draw[oob] (held.north east) -- (t2.south west);
\draw[oob] (held.east) -- (t3.south west);
\draw[oob] (t2.south east) -- (oobavg.north west);
\draw[oob] (t3.east) -- (oobavg.west);
\end{tikzpicture}

\caption{\textcolor{cyan!65!black}{Blue} bootstraps train \textcolor{green!45!black}{green} trees with $m$ split candidates; black arrows aggregate, \textcolor{red!70!black}{red dashed} arrows give OOB prediction for $i$.}
\label{fig:sl-random-forest-oob-decorrelation}
\end{figure}

\begin{definition} \label{def:sl-cart-regression-tree-scores} \hlt{CART Regression-Tree Notation and Scores}\\
Let $D_n=\{(x_i,y_i)\}_{i=1}^n\subseteq\mathcal{X}\times\R$. For $R\subseteq\mathcal{X}$, set $I_R\defeq\{i:x_i\in R\}$ and $n_R\defeq |I_R|$.
\begin{enumerate}[label=(\roman*)]
\setlength{\itemsep}{0pt}
\item \hlt{Cell loss}: if $n_R>0$, define $\bar y_R\defeq n_R^{-1}\sum\limits_{i\in I_R}y_i$ and $\mathrm{SSE}(R)\defeq\sum\limits_{i\in I_R}(y_i-\bar y_R)^2$.
\item \hlt{Tree predictor}: if $T$ has leaves $R_1(T),\ldots,R_{|T|}(T)$, let $m_T(x)$ be the unique $m\in\{1,\ldots,|T|\}$\\
with $x\in R_m(T)$ and set $\widehat{f}_T(x)\defeq\bar y_{R_{m_T(x)}(T)}$.
\item \hlt{Split children}: for $a=(j,s)\in\{1,\ldots,p\}\times\R$, set $R_L(a)\defeq R\cap\{x:x_j\leq s\}$ and $R_R(a)\defeq R\cap\{x:x_j>s\}$.
\item \hlt{Admissible splits}: set $\mathcal{A}_R\defeq\{a\in\{1,\ldots,p\}\times\R:n_{R_L(a)}>0,n_{R_R(a)}>0\}$.
\item \hlt{Split score}: for $a\in\mathcal{A}_R$, set $S_R(a)\defeq\mathrm{SSE}(R_L(a))+\mathrm{SSE}(R_R(a))$.
\item \hlt{Pruning score}: for $\alpha\geq0$, set $C_\alpha(T)\defeq\sum\limits_{m=1}^{|T|}\mathrm{SSE}(R_m(T))+\alpha|T|$.
\end{enumerate}
\end{definition}

\begin{method} \label{meth:sl-cart-tree-growing-and-pruning} \hlt{CART-Style Regression Trees}
\begin{enumerate}[label=(\roman*)]
\setlength{\itemsep}{0pt}
\item \hlt{Grow}: split eligible leaves by $a_R\in\operatorname*{argmin}\limits_{a\in\mathcal{A}_R}S_R(a)$ until stopping, yielding $T_0$.
\item \hlt{Prune}: for each $\alpha\geq0$, choose $T_\alpha\in\operatorname*{argmin}\limits_{T\preceq T_0}C_\alpha(T)$.
\item \hlt{Select}: choose $\widehat{\alpha}\geq0$ by validation or cross-validation and return $T_{\widehat{\alpha}}$.
\end{enumerate}
Deep growth lowers bias; the cost-complexity penalty $\alpha|T|$ controls variance by subtree selection.
\end{method}

\begin{definition} \label{def:sl-classification-tree-impurity} \hlt{Classification Tree Impurity}\\
For node $m$, let $N_{mk}\defeq|\{i:x_i\in R_m,y_i=k\}|$, $N_m\defeq\sum\limits_{k=1}^K N_{mk}$, $\widehat{p}_{mk}\defeq \frac{N_{mk}}{N_m}$. Common impurity scores are:
\begin{enumerate}[label=(\roman*)]
\setlength{\itemsep}{0pt}
\item \hlt{Misclassification}: $I_{\mathrm{mis}}(m)\defeq1-\max_{1\leq k\leq K}\widehat{p}_{mk}$.
\item \hlt{Gini}: $I_{\mathrm{Gini}}(m)\defeq\sum_{k=1}^K\widehat{p}_{mk}(1-\widehat{p}_{mk})=1-\sum_{k=1}^K\widehat{p}_{mk}^2$.
\item \hlt{Entropy}: $I_{\mathrm{ent}}(m)\defeq-\sum_{k:\widehat{p}_{mk}>0}\widehat{p}_{mk}\log\widehat{p}_{mk}$.
\end{enumerate}
The usual leaf prediction is the majority class, $\widehat{k}(m)\in\operatorname*{argmax}\limits_{1\leq k\leq K}\widehat{p}_{mk}$. \\
With unequal losses $L_{kk}=0$, choose the class minimizing expected loss, $\widehat{\ell}(m)\in\operatorname*{argmin}\limits_{1\leq\ell\leq K}\sum_{k=1}^K L_{k\ell}\widehat{p}_{mk}$.
\end{definition}

\begin{remark} \label{rem:sl-tree-practical-issues} \hlt{Tree Practicalities}\\
For categorical levels $\mathcal{A}_j$, split $B\subset\mathcal{A}_j$ versus $B^c$.\\
For binary responses, set $I_a\defeq\{i:X_{ij}=a\}$ and order levels by $\widehat{p}_a\defeq |I_a|^{-1}\sum_{i\in I_a}Y_i$.\\
Missing values can be an $\mathrm{NA}$ level or use a surrogate split $g'$ maximising $\sum_i\mathbf{1}\{g'(x_i)=g(x_i)\}$.\\
Many candidate splits inflate gain; control growth with $N_m\geq n_{\min}$, depth limits, validation, or pruning.
\end{remark}

\begin{method} \label{meth:sl-bagging-prediction-aggregation} \hlt{Bagged Prediction and Votes}\\
Let $D_n\defeq\{(x_i,y_i)\}_{i=1}^n$ and $B\geq1$; base outputs are $\widehat{f}$ for regression and $(\widehat{g},\widehat{p}_1,\ldots,\widehat{p}_K)$ for classification.
\begin{enumerate}[label=(\roman*)]
\setlength{\itemsep}{0pt}
\item For $b=1,\ldots,B$, draw a bootstrap sample $D_b^*$ from $D_n$ and fit the base learner, giving $\widehat{f}^{*b}$ or $(\widehat{g}^{*b},\widehat{p}_{b1},\ldots,\widehat{p}_{bK})$.
\item Return the task-specific aggregate: for regression, $\widehat{f}_{\mathrm{bag}}(x)\defeq B^{-1}\sum\limits_{b=1}^B\widehat{f}^{*b}(x)$; for classification, either $\widehat{g}_{\mathrm{vote}}(x)\in\operatorname*{argmax}\limits_{1\leq k\leq K}\sum\limits_{b=1}^B\mathbf{1}\{\widehat{g}^{*b}(x)=k\}$ or $\bar{p}_k(x)\defeq B^{-1}\sum\limits_{b=1}^B\widehat{p}_{bk}(x)$ with $\widehat{g}_{\mathrm{prob}}(x)\in\operatorname*{argmax}\limits_{1\leq k\leq K}\bar{p}_k(x)$.
\end{enumerate}
Averaging reduces perturbation variance for unstable base learners.\\
Random split-feature subsets additionally reduce $\operatorname{corr}(T_b,T_{b'})$.
\end{method}

\begin{theorem} \label{thm:sl-bagging-squared-error-aggregation} \hlt{Squared-Error Benefit of Averaging}\\
Fix $x$, let $Z\defeq\widehat{f}^{\ast}(x)$ be perturbation-random and independent of $Y\mid X=x$, and set $f_{\mathrm{ag}}(x)\defeq\mathbb{E}_{\ast}Z$.\\
If $\mathbb{E}_{\ast}Z^2<\infty$ and $\mathbb{E}_Y Y^2<\infty$, then $\mathbb{E}_{\ast,Y}[(Y-Z)^2]=\mathbb{E}_Y[(Y-f_{\mathrm{ag}}(x))^2]+\mathbb{E}_{\ast}[(Z-f_{\mathrm{ag}}(x))^2]$.
\begin{proof} Let $\mu\defeq f_{\mathrm{ag}}(x)=\mathbb{E}_{\ast}Z$.\\
Since $Y-\mu$ and $Z-\mu$ are independent and $\mathbb{E}_{\ast}(Z-\mu)=0$, $\mathbb{E}_{\ast,Y}[(Y-\mu)(Z-\mu)]=\mathbb{E}_Y(Y-\mu)\mathbb{E}_{\ast}(Z-\mu)=0$.\\
Expanding $Y-Z=(Y-\mu)-(Z-\mu)$ gives $\mathbb{E}_{\ast,Y}[(Y-Z)^2]=\mathbb{E}_Y[(Y-\mu)^2]+\mathbb{E}_{\ast}[(Z-\mu)^2]$.\\
For independent fits $Z_b$, the average $\widehat{f}_B(x)\defeq B^{-1}\sum\limits_{b=1}^B Z_b$ has mean $\mu$ and perturbation variance $\operatorname{Var}_{\ast}(Z)/B$.\\
\hlt{Structural Commentary}: squared loss decomposes into target risk plus perturbation variance.\\
\hlt{Key Insight}: bagging estimates $\mathbb{E}_{\ast}\widehat{f}^{\ast}(x)$; unstable base learners gain most because $\operatorname{Var}_{\ast}(Z)$ is large.
\end{proof}
\end{theorem}


\begin{theorem} \label{thm:sl-correlated-ensemble-variance} \hlt{Correlated Ensemble Variance}\\
Let $Z_1,\ldots,Z_B$ have $\operatorname{Var}(Z_b)=\sigma^2>0$ and $\operatorname{Corr}(Z_b,Z_c)=\rho$ for $b\neq c$.\\
For $\bar{Z}_B\defeq \frac{1}{B}\sum\limits_{b=1}^B Z_b$, $\operatorname{Var}(\bar{Z}_B)=\rho\sigma^2+\frac{1-\rho}{B}\sigma^2$; hence, if $\rho\geq0$, $\operatorname{Var}(\bar{Z}_B)\to\rho\sigma^2$ as $B\to\infty$.
\begin{proof}
Since $\operatorname{Cov}(Z_b,Z_c)=\rho\sigma^2$ for $b\neq c$, $\operatorname{Var}(\bar{Z}_B)=B^{-2}\left\{\sum\limits_{b=1}^B\operatorname{Var}(Z_b)+2\sum\limits_{1\leq b<c\leq B}\operatorname{Cov}(Z_b,Z_c)\right\}$.\\
Substitution gives $\operatorname{Var}(\bar{Z}_B)=B^{-2}\left\{B\sigma^2+B(B-1)\rho\sigma^2\right\}=\rho\sigma^2+\frac{1-\rho}{B}\sigma^2$.\\
\hlt{Structural Commentary}: $\operatorname{Var}(\bar{Z}_B)=\rho\sigma^2+(1-\rho)\sigma^2/B$ separates shared and idiosyncratic variance.\\
\hlt{Key Insight}: increasing $B$ shrinks only $(1-\rho)\sigma^2/B$; Random Forest lowers $\rho$ by split randomisation.
\end{proof}
\end{theorem}

\begin{method} \label{meth:sl-random-forest-oob-and-importance} \hlt{Out-of-Bag Error and Forest Importance}\\
Let $I_b^*$ index bootstrap sample for $T_b$, $O_i\defeq\{b:i\notin I_b^*\}$, $\mathcal{I}_{\mathrm{oob}}\defeq\{i:|O_i|>0\}$, and $L$ be prediction loss.
\begin{enumerate}[label=(\roman*)]
\setlength{\itemsep}{0pt}
\item For each $i\in\mathcal{I}_{\mathrm{oob}}$, construct $\widehat{h}_{\mathrm{oob}}(x_i)$: for regression, set $\widehat{h}_{\mathrm{oob}}(x_i)\defeq |O_i|^{-1}\sum\limits_{b\in O_i}T_b(x_i)$; for classification, use an OOB vote or averaged class-probability vector over $b\in O_i$.
\item Set $\widehat{R}_{\mathrm{oob}}\defeq |\mathcal{I}_{\mathrm{oob}}|^{-1}\sum\limits_{i\in\mathcal{I}_{\mathrm{oob}}}L(y_i,\widehat{h}_{\mathrm{oob}}(x_i))$ and update the estimate as $B$ changes.
\end{enumerate}
For feature $j$, let $\mathcal{S}_{bj}$ be $j$-splits of $T_b$, with gains $\Delta Q_{bs}$ and OOB losses $R_b^{\mathrm{oob}},R_{bj}^{\pi}$ before/after permuting $X_j$.\\
Define $I_j^{\mathrm{split}}\defeq\sum\limits_{b=1}^B\sum\limits_{s\in\mathcal{S}_{bj}}\Delta Q_{bs}$ and $I_j^{\pi}\defeq B^{-1}\sum\limits_{b=1}^B(R_{bj}^{\pi}-R_b^{\mathrm{oob}})$.\\
Both are fitted-forest summaries; $I_j^{\pi}$ is not a refit-without-$X_j$ loss contrast.
\end{method}

\begin{remark} \label{rem:sl-random-forest-proximity-overfitting} \hlt{Proximities, Tuning, and Overfitting}\\
Let $L_b(i)$ denote the terminal leaf of $x_i$ in $T_b$, and set $O_{ij}\defeq\{b:i,j\notin I_b^*\}$.\\
Define $P_{ij}\defeq B^{-1}\sum\limits_{b=1}^B\mathbf{1}\{L_b(i)=L_b(j)\}$ and, when $|O_{ij}|>0$, $P_{ij}^{\mathrm{oob}}\defeq |O_{ij}|^{-1}\sum\limits_{b\in O_{ij}}\mathbf{1}\{L_b(i)=L_b(j)\}$.\\
Thus $P$ is a fitted-forest similarity; it is neither data-intrinsic nor generally a metric.\\
Tuning effects are parameter-specific: $m\downarrow$ lowers $\operatorname{corr}(T_b,T_{b'})$ and can raise bias; $n_{\min}\uparrow$ smooths terminal leaves; $B\uparrow$ lowers finite-ensemble Monte Carlo error but does not regularise the limiting forest.\\
Defaults $m_{\mathrm{cls}}\approx\sqrt p$ and $m_{\mathrm{reg}}\approx p/3$ are starting values only.
\end{remark}

\begin{method} \label{meth:sl-error-correcting-output-codes} \hlt{Error-Correcting Output Codes}\\
For training data $(x_i,y_i)_{i=1}^n$ with $y_i\in\{1,\ldots,K\}$, let $d_H$ denote Hamming distance.\\
A code matrix $C\in\{0,1\}^{K\times L}$ assigns class $k$ the row $C_{k\cdot}$, and column $\ell$ induces binary labels $z_{i\ell}\defeq C_{y_i\ell}$.
\begin{enumerate}[label=(\roman*)]
\setlength{\itemsep}{0pt}
\item Choose $C$ with distinct rows, equivalently $d_H(C_{k\cdot},C_{k'\cdot})>0$ for $k\neq k'$.
\item For each $\ell=1,\ldots,L$, fit $\widehat{p}_\ell(x)\approx\mathbb{P}(Z_\ell=1\mid X=x)$ to $(x_i,z_{i\ell})_{i=1}^n$.
\item For a test point $x$, compute $\delta_k(x)\defeq\sum\limits_{\ell=1}^L |C_{k\ell}-\widehat{p}_\ell(x)|$ for $k=1,\ldots,K$ and return $\widehat{k}(x)\in\operatorname*{argmin}\limits_{1\leq k\leq K}\delta_k(x)$.
\end{enumerate}
Row separation gives error correction; redundant columns reduce variance and absorb column errors.
\end{method}

\begin{method} \label{meth:sl-stacking-and-bumping} \hlt{Stacking and Bumping}\\
Let $\widehat{f}_1,\ldots,\widehat{f}_M$ be fitted candidates, $\widehat{f}^{-i}_m(x_i)$ the out-of-sample prediction for observation $i$ from candidate $m$, $\mathcal{W}\subseteq\R^M$ the admissible weight set, and $L$ a loss. Commonly $\mathcal{W}\defeq\{w\in\R^M:w_m\geq0,\sum\limits_{m=1}^M w_m=1\}$.
\begin{enumerate}[label=(\roman*)]
\setlength{\itemsep}{0pt}
\item \hlt{Stacking}: set $R_{\mathrm{st}}(w)\defeq\sum\limits_{i=1}^n(y_i-\sum\limits_{m=1}^M w_m\widehat{f}^{-i}_m(x_i))^2$.\\
Choose $\widehat{w}^{\mathrm{st}}\in\operatorname*{argmin}\limits_{w\in\mathcal{W}}R_{\mathrm{st}}(w)$, and return $\widehat{f}_{\mathrm{st}}(x)\defeq\sum\limits_{m=1}^M\widehat{w}^{\mathrm{st}}_m\widehat{f}_m(x)$.
\item \hlt{Bumping}: fit bootstrap candidates $\widehat{f}^{*1},\ldots,\widehat{f}^{*B}$, set $R_b^*\defeq\sum\limits_{i=1}^n L(y_i,\widehat{f}^{*b}(x_i))$.\\
Choose $\widehat{b}\in\operatorname*{argmin}\limits_{1\leq b\leq B}R_b^*$, and return $\widehat{f}^{*\widehat{b}}$.
\end{enumerate}
Coordinate weights $w=e_m$ give cross-validated model selection; general $w\in\mathcal{W}$ gives model averaging.\\
Stacking averages fixed candidates; bumping selects among resampled fits.
\end{method}
